\section{Inherited weighted quartet analytic estimate}
\label{sec:quartet-upper}


The following proposition, including its constants and proof strategy,
is the finalized Phase~9 weighted quartet theorem. It is reproduced here
so the upper-bound argument can be read without another file. It is not
a new Phase~10 analytic estimate. The source is the weighted-quartet theorem in the Phase~9 working
paper~\cite{huang9}; the exact source path and reviewed hash are recorded
in the accompanying build receipt.

\begin{proposition}[Phase 9 weighted quartet estimate]
\label{prop:inherited}
For four distinct weighted states and nonnegative nonzero polynomial
curvature of degree at most $r\ge1$, let $R$ be the unrestricted
deterministic hierarchy price at capacities $(2,3)$. If $R\ge128$, then
\[
 R^2\le16r^2e^{20r/\sqrt R}.
\]
If $R\ge256$, then $R(\log R)^2\le1936r^2$. In particular, for $r\ge2$,
$R\le1936r^2/(\log r)^2$.
\end{proposition}

\begin{proof}
We record two classical polynomial estimates. If $f\ge0$ has degree at
most $r$ on an interval $E$ of length $L$, and $M_E=\max_Ef$, the classical Markov
inequality~\cite{shadrin} $\|f'\|_\infty\le2r^2M_E/L$ implies
\begin{equation}
 \int_E f\ge\frac{LM_E}{8r^2}.
 \label{eq:markov}
\end{equation}
Indeed a one-sided segment of length $L/(4r^2)$ is available from a
maximizer, and $f\ge M_E/2$ there. If $t_0$ lies at distance $d>0$ to
the right of $E$, Chebyshev extremality gives
\begin{equation}
 |f(t_0)|\le M_E T_r(1+2d/L)
 \le M_Ee^{2r\sqrt{d/L}}.
 \label{eq:cheb}
\end{equation}
For completeness, after rescaling $E$ to $[-1,1]$, any violation of the
first inequality would yield $\lambda\in(0,1)$ such that
$T_r-\lambda f/M_E$ has $r$ roots between the alternating extrema of
$T_r$ and an additional root at the external point. This contradicts
its degree. The exponential bound follows from
$T_r(\cosh s)=\cosh(rs)$ and
$\operatorname{arcosh}(1+x)\le\sqrt{2x}$.

Since $R>1$, the argument of Theorem~\ref{thm:reduction}, specialized to
four states, permits the notation
\[
 A<B<C<D,\quad q=D(BC)=\OPT_3,\quad
 Q=D(AB)+D(CD)=\OPT_2.
\]
Write $p_L=D(AB)$, $p_R=D(CD)$, $T_L=D(ABC)$, and $T_R=D(BCD)$.
Three feasible deterministic hierarchies give
\begin{equation}
 R\le\frac{\min(p_L,p_R)}q,\qquad
 R\le\frac{T_L}Q,\qquad R\le\frac{T_R}Q.
 \label{eq:candidates}
\end{equation}
For the first retain $AB\mid CD$ and merge the cheaper outer pair at
the fine level. For the others retain fine pair $BC$ and merge the
corresponding adjacent triple at the coarse level. All right sides
are at least one by the defining flat optimality. No enumeration of
noncontiguous optimal hierarchies is needed.

Reflect the source if needed so that the source weights satisfy
$p_B\le p_C$. Normalize only the coordinate of the triple $BCD$ to
\[
 B=0,\qquad C=c\in(0,1),\qquad D=1,
 \qquad (p_B,p_C,p_D)=(u,v,w),\quad 0<u\le v.
\]
The transformed curvature is $L_0^2g(B+L_0t)$, where $L_0$ is the
original triple hull length; its degree and nonnegativity are preserved.
Weights are unchanged. Denote that curvature again by $g$. The triple,
adjacent-pair, and outer-pair kernels are
\begin{align*}
 K(t)&=\min\{ut+v(t-c)_+,\ w(1-t)+v(c-t)_+\},\\
 k_q(t)&=\min\{ut,v(c-t)\}\mathbf1_{[0,c]}(t),\\
 k_p(t)&=\min\{v(t-c),w(1-t)\}\mathbf1_{[c,1]}(t),\\
 H(t)&=\min\{ut,w(1-t)\}.
\end{align*}
Put $T=T_R=\int Kg$ and $S=p_R+q=\int(k_q+k_p)g$. From
\eqref{eq:candidates},
\begin{equation}
 q\le\frac{Q}{2R}\le\frac{T}{2R^2},\qquad
 S\le Q+q\le\frac{2T}{R}.
 \label{eq:small}
\end{equation}
Two pointwise kernel comparisons control degenerating weights:
\begin{equation}
 0\le K-H\le k_q+k_p,
 \label{eq:H}
\end{equation}
and, in their respective open intervals,
\begin{equation}
 \frac{k_q(t)}{K(t)}\ge\min\{1,(c-t)/t\}\quad(t<c),
 \qquad
 \frac{k_p(t)}{K(t)}\ge\frac{t-c}{c+2(t-c)}\quad(t>c).
 \label{eq:ratios}
\end{equation}
For $t<c$, subtracting $H$ from $K$ amounts to adding $v(c-t)$ to one
entry of a minimum, so the increment is between zero and
$\min\{ut,v(c-t)\}=k_q$. For $t>c$, adding $v(t-c)$ to the other entry
gives the bound by $k_p$. For the first ratio, use $K\le ut$ and
$v\ge u$. For the second, write $\Delta=t-c>0$ and use
\[
 \frac{\min\{B,C\}}{\min\{A+B,C\}}\ge\frac{B}{A+B},
 \quad A=ut,\ B=v\Delta,\ C=w(1-t),
\]
followed by $v\ge u$. Endpoint claims follow by continuity.

Assume $R\ge128$ and set
$\varepsilon=16/R\le1/8$ and
$J=\{t\in[0,1]:|t-c|\le\varepsilon c\}$. Outside $J$,
\eqref{eq:ratios} yields
$k_q+k_p\ge\varepsilon K/(1+2\varepsilon)$. Consequently,
\[
 \int_{[0,1]\setminus J}Kg
 \le\frac{1+2\varepsilon}{\varepsilon}S
 \le\frac{5T}{32},\qquad
 \int_JHg\ge T-\frac{5T}{32}-S\ge\frac{3T}{4}.
\]
Here $S\le2T/R\le T/64$. Since
$H(t)\le u(1+\varepsilon)c$ on $J$ and $|J|\le2\varepsilon c$,
there is $t_0\in J$ such that
\[
 g(t_0)\ge\frac{T}{4\varepsilon u c^2}.
\]
Consider $E=[c/4,c(1-2\varepsilon)]$. It has length at least $c/2$;
$t_0$ is to its right at distance at most $3\varepsilon c$; and
$k_q\ge2u\varepsilon c$ on $E$. Using \eqref{eq:cheb},
\[
 M_E\ge\frac{T}{4\varepsilon u c^2}
 e^{-\sqrt{24}\,r\sqrt\varepsilon}
 \ge\frac{T}{4\varepsilon u c^2}e^{-20r/\sqrt R}.
\]
Equation~\eqref{eq:markov} now implies
\[
 q\ge2u\varepsilon c\int_Eg
 \ge\frac{u\varepsilon c^2M_E}{8r^2}
 \ge\frac{T}{32r^2}e^{-20r/\sqrt R}.
\]
Together with \eqref{eq:small}, this proves the implicit inequality.

Set $z=r/\sqrt R$. The inequality implies
\[
 \log(R/16)\le20z+2\log z\le22z.
\]
For $R\ge256$, the left side is at least $\tfrac12\log R$, yielding
$\sqrt R\log R\le44r$, hence $R(\log R)^2\le1936r^2$.
For $r\ge2$ and $R\ge256$ with $R>r$, use $\log R\ge\log r$.
If $R\le r$, use $(\log r)^2\le r$; if $R<256$, use
$(\log r)^2\le r^2$ and $256<1936$. These cases give the explicit bound.
\end{proof}

